\item\;  解答は特別な指示がない限り，結果だけでなく過程を含め記述すること
\item\;  記述の際に使用してよい筆記具は鉛筆，シャープペンシル，ボールペンのいずれかである．
\item\;  試験開始の合図があるまで問題用紙を開かないこと
\item\;  解答用紙及び問題用紙の所定の欄に受験番号を記述すること
\item\;  問題は\,\fbox{1}\,から\,\fbox{\amountq}\,まである
\item\;  \fbox{\hreq}\,から\,\fbox{\ereq}\,までは必答問題，\fbox{\hselect}\,から\,\fbox{\eselect}\,までは選択問題であり，選択問題は$\amountsel$問選択して答え\\[2mm]
         \hspace{0.2cm}ること．
\item\;  問題冊子は各自持ち帰ること
\item\;  体調不良等を除き試験中の途中退室は認められない
\item\;  誤字・脱字や落丁，問題の不備等が有る場合は静かに手を上げて試験官に知らせること
\item\;  携帯電話等の電子機器は電源を切りカバン等にしまうこと．
\item\;  問題の一部及びすべてを許可なく複製，またはSNS等のインターネット上に掲載することは\\[2mm]
         \hspace{0.2cm}一切禁止である．